\chapter{Podrozmaitości przestrzeni euklidesowej}
\label{ma2:ch:submanifolds}

Rozważamy podrozmaitości bez brzegu klasy $C^s$,
gdzie $s$ jest dodatnią liczbą całkowitą lub $s=\infty$.
Ogólną teorię rozmaitości rozwinięto w DG-01; tutaj
koncentrujemy się na opisie w przestrzeni euklidesowej
i rachunkach potrzebnych do ekstremów związanych oraz całkowania.

\section{Trzy lokalne opisy podrozmaitości}

\begin{definicja}
\label{ma2:def:submanifold}
Podzbiór $S\subseteq\R^m$ jest podrozmaitością
zanurzoną wymiaru $k$, $0\le k\le m$, klasy $C^s$,
jeżeli dla każdego $p\in S$ istnieją otwarte
otoczenie $U\subseteq\R^m$ punktu $p$ oraz dyfeomorfizm
$\alpha:U\to W\subseteq\R^k\times\R^{m-k}$
klasy $C^s$ na zbiór otwarty $W$, taki że
$\alpha(p)=0$ i
\[
 \alpha(S\cap U)=W\cap(\R^k\times\{0\}).
\]
Taką mapę nazywamy mapą plasterkową.
\end{definicja}

Topologię $S$ stanowi topologia podprzestrzeni;
jest ona Hausdorffa i ma przeliczalną bazę odziedziczoną
z $\R^m$. Oznaczamy przez $\pi_1,\pi_2$ rzuty na
pierwszą i drugą składową produktu.
Mapa $\pi_1\circ\alpha|_{S\cap U}$ ma otwarty obraz
$\{u:(u,0)\in W\}$, ponieważ jest to przeciwobraz
$W$ przez ciągłe włożenie $u\mapsto(u,0)$.
Jej odwrotność $u\mapsto\alpha^{-1}(u,0)$ jest ciągła.
Otrzymujemy zatem homeomorfizmy, których przejścia są klasy
$C^s$: mają postać $u\mapsto\pi_1(\beta(\alpha^{-1}(u,0)))$
dla dwóch map plasterkowych $\alpha,\beta$.
Dziedziną przejścia jest obraz przecięcia dziedzin
obu map na $S$ przez pierwszą z nich; jest to zbiór
otwarty w $\R^k$, ponieważ obie dziedziny są względnie
otwarte w $S$, a pierwsza mapa jest homeomorfizmem.
Odwrotne przejście otrzymujemy przez zamianę map.
Dla $s=\infty$ pełna konstrukcja struktury rozmaitości
znajduje się w
\atlasref{DG-01::stw:podrozm-struktura}{DG-01: struktura podrozmaitości}.

\begin{twierdzenie}[Lokalne równania i wykresy]
\label{ma2:thm:submanifold-descriptions}
Ustalmy $0\le k\le m$, $S\subseteq\R^m$ i $p\in S$.
Następujące warunki są równoważne:
\begin{enumerate}
\item $S$ ma w $p$ mapę plasterkową rzędu $k$ klasy $C^s$.
\item Istnieje otwarte otoczenie $U$ punktu $p$
i odwzorowanie $F:U\to\R^{m-k}$ klasy $C^s$ takie, że
$S\cap U=F^{-1}(\{0\})$ oraz $DF(p)$ jest surjektywny.
\item Po permutacji współrzędnych, na pewnym otoczeniu
produktowym $A\times B$ punktu $p$ zachodzi
$S\cap(A\times B)=\{(x,\varphi(x)):x\in A\}$,
gdzie $\varphi:A\to B$ jest klasy $C^s$, a
$A\subseteq\R^k$, $B\subseteq\R^{m-k}$ są otwarte.
\end{enumerate}
\end{twierdzenie}
\begin{proof}
Dla pierwszego warunku wybieramy $F=\pi_2\circ\alpha$.
Równość zbiorów wynika z definicji mapy plasterkowej:
$F(q)=0$ wtedy i tylko wtedy, gdy
$\alpha(q)\in W\cap(\R^k\times\{0\})$, czyli $q\in S\cap U$.
Z reguły łańcuchowej (twierdzenie \ref{ma2:thm:frechet-chain})
mamy $DF(p)=\pi_2D\alpha(p)$. Dla każdego $w\in\R^{m-k}$
wektor $(D\alpha(p))^{-1}(0,w)$ przechodzi przez $DF(p)$
na $w$, więc operator jest surjektywny.

Z drugiego warunku, jeśli $m-k>0$, surjektywność
$DF(p)$ daje rząd $m-k$. Wybieramy więc $m-k$
liniowo niezależnych kolumn jego macierzy, tworzących
kwadratową podmacierz o niezerowym wyznaczniku. Przenosimy te współrzędne na koniec
i piszemy $p=(a,b)$. Wtedy $D_yF(a,b)$ jest
izomorfizmem $\R^{m-k}\to\R^{m-k}$.
Ponieważ $p\in S\cap U$, mamy $F(a,b)=0$.
Twierdzenie o funkcji uwikłanej
(\ref{ma2:thm:implicit-function}) daje otoczenia
$A,B$ z $A\times B\subseteq U$ oraz funkcję
$\varphi:A\to B$ spełniającą
$F(x,y)=0$ wtedy i tylko wtedy, gdy $y=\varphi(x)$.
Ponieważ $S\cap U=F^{-1}(\{0\})$, daje to dokładnie
żądaną równość na produkcie, a regularność funkcji
$\varphi$ wynika z tego samego twierdzenia.
Gdy $m-k=0$, przeciwdziedziną $F$ jest $\{0\}$,
więc $S\cap U=U$; wystarczy wybrać otwarte
otoczenie $A\subseteq U$ punktu $p$, $B=\{0\}$
oraz $\varphi(x)=0$. W przypadku $k=0<m$
twierdzenie o funkcji uwikłanej daje $A=\{0\}$,
więc $S\cap(A\times B)$ składa się tylko z $p$.
Podrozmaitości wymiaru zero są zatem zbiorami
dyskretnymi w topologii podprzestrzeni, a podrozmaitości
wymiaru $m$ są otwartymi podzbiorami $\R^m$.

Z trzeciego warunku, przy $p=(a,\varphi(a))$,
definiujemy $\alpha(x,y)=(x-a,y-\varphi(x))$.
Na obrazie odwrotność ma wzór
$(u,v)\mapsto(u+a,v+\varphi(u+a))$.
Obraz jest otwarty: składa się z tych $(u,v)$,
dla których $u+a\in A$ i $v+\varphi(u+a)\in B$,
czyli jest przeciwobrazem otwartego produktu przez
ciągłe odwzorowanie określone na $(A-a)\times\R^{m-k}$.
Oba wzory są klasy $C^s$. Istotnie, podstawienie
$(u,v)=(x-a,y-\varphi(x))$ do wzoru odwrotnego daje
$(x,y)$, a podstawienie $(x,y)=(u+a,v+\varphi(u+a))$
do wzoru na $\alpha$ daje $(u,v)$.
Współrzędne po ewentualnej permutacji identyfikujemy
z pierwotnymi przez odwracalną mapę liniową.
Ponadto $\alpha(p)=0$ oraz
druga składowa $\alpha$ jest zerowa dokładnie na wykresie.
Otrzymujemy mapę plasterkową.
\end{proof}

\section{Parametryzacje i przestrzenie styczne}

\begin{definicja}
\label{ma2:def:local-parametrization}
Niech $S$ będzie podrozmaitością wymiaru $k$
klasy $C^s$. Parametryzacją lokalną $S$ nazywamy odwzorowanie
$\sigma:V\to S$, gdzie $V\subseteq\R^k$ jest zbiorem otwartym,
które jest klasy $C^s$ jako odwzorowanie do $\R^m$,
jest homeomorfizmem na względnie otwarty podzbiór $S$
i ma różnowartościową pochodną w każdym punkcie $V$.
\end{definicja}

Z mapy plasterkowej otrzymujemy parametryzację
$\sigma(u)=\alpha^{-1}(u,0)$ na otwartym zbiorze
$V=\{u:(u,0)\in W\}$. Ma ona obraz $S\cap U$,
gdyż $\alpha(S\cap U)=W\cap(\R^k\times\{0\})$.
Jej odwrotność to $\pi_1\circ\alpha$ ograniczone do
$S\cap U$; oba odwzorowania są ciągłe.
Pochodna jest złożeniem włożenia $h\mapsto(h,0)$
z izomorfizmem $D\alpha^{-1}(u,0)$, więc jest
różnowartościowa. W szczególności wykres ma
parametryzację $\sigma(x)=(x,\varphi(x))$.

\begin{definicja}
\label{ma2:def:tangent-submanifold}
Niech $p\in S$, gdzie $S$ jest podrozmaitością.
Przestrzeń styczna $T_pS$ to zbiór prędkości
$\gamma'(0)$ krzywych $\gamma:(-\varepsilon,\varepsilon)\to S$
klasy $C^1$ (jako odwzorowania do $\R^m$), dla których
$\gamma(0)=p$. Przestrzeń normalną definiujemy jako
$N_pS=(T_pS)^\perp$ względem standardowego iloczynu skalarnego.
Na tym etapie $T_pS$ jest zbiorem; jego liniowość
wykażemy w następnym twierdzeniu. Afiniczną przestrzenią
styczną w $p$ nazywamy $p+T_pS$, a normalną $p+N_pS$.
\end{definicja}

\begin{twierdzenie}
\label{ma2:thm:tangent-descriptions}
Dla podrozmaitości $S$ wymiaru $k$, parametryzacji
lokalnej $\sigma$ z $\sigma(a)=p$ oraz lokalnego
opisu równaniami $F=0$ z twierdzenia
\ref{ma2:thm:submanifold-descriptions} zachodzi
\[
 T_pS=D\sigma(a)(\R^k)=\ker DF(p).
\]
W szczególności $\dim T_pS=k$. Jeżeli
$F=(F_1,\ldots,F_{m-k})$, to
\[
 N_pS=\operatorname{span}\{\nabla F_j(p):1\le j\le m-k\},
 \qquad \R^m=T_pS\oplus N_pS.
\]
\end{twierdzenie}
\begin{proof}
Dla każdej krzywej z definicji przestrzeni stycznej
ciągłość zapewnia, że w otoczeniu zera $\gamma(t)\in U$,
gdzie $S\cap U=F^{-1}(\{0\})$. Różniczkując
$F(\gamma(t))=0$, otrzymujemy $DF(p)\gamma'(0)=0$.
Zatem $T_pS\subseteq\ker DF(p)$.
Dla dowolnego $h\in\R^k$ krzywa
$\gamma_h(t)=\sigma(a+th)$ jest określona dla małych
$|t|$, leży w $S$ i ma prędkość $D\sigma(a)h$.
W konsekwencji
\[
 D\sigma(a)(\R^k)\subseteq T_pS\subseteq\ker DF(p).
\]
Skrajne podprzestrzenie mają wymiar $k$: pierwsza
z różnowartościowości $D\sigma(a)$, druga z surjektywności
$DF(p)$ i twierdzenia o wymiarze jądra i obrazu.
Obie inkluzje muszą być równościami.
Dalej, z definicji gradientu w $\R^m$,
\[
 DF(p)v=(\langle\nabla F_1(p),v\rangle,\ldots,
              \langle\nabla F_{m-k}(p),v\rangle).
\]
Zatem $v\in\ker DF(p)$ dokładnie wtedy, gdy
$\langle\nabla F_j(p),v\rangle=0$ dla każdego $j$.
Biorąc dopełnienia ortogonalne i korzystając ze
skończonego wymiaru, otrzymujemy wzór na $N_pS$.
Rozkład ortogonalny daje ostatnią równość.
Dla $k=m$ pusta rodzina gradientów rozpina $\{0\}$,
a dla $k=0$ przestrzeń styczna jest zerowa.
Dla $k<m$ gradienty są liniowo niezależne, ponieważ
są wierszami macierzy surjektywnego operatora $DF(p)$
(utożsamionymi z wektorami przez iloczyn skalarny).
Tworzą więc bazę przestrzeni normalnej.
\end{proof}

W wersji gładkiej są to odpowiedniki wyników
\atlasref{DG-01::stw:styczna-podrozmaitosci}{DG-01: przestrzeń styczna}
oraz
\atlasref{DG-01::stw:styczna-normalne-euklidesowe}{DG-01: wektory normalne}.

\begin{wniosek}
\label{ma2:cor:graph-tangent}
Niech $A\subseteq\R^k$ będzie otwarty,
$\varphi:A\to\R^{m-k}$ będzie klasy $C^s$ i
$S=\{(x,\varphi(x)):x\in A\}$.
Dla $a\in A$ oraz $p=(a,\varphi(a))$ mamy
\[
 T_pS=\{(h,D\varphi(a)h):h\in\R^k\},\qquad
 N_pS=\{(-D\varphi(a)^*w,w):w\in\R^{m-k}\}.
\]
\end{wniosek}
\begin{proof}
Niech $D\varphi(a)^*$ oznacza operator sprzężony
względem euklidesowych iloczynów skalarnych.
Z pochodnej parametryzacji $\sigma(x)=(x,\varphi(x))$
otrzymujemy $D\sigma(a)h=(h,D\varphi(a)h)$,
więc pierwszy wzór wynika z twierdzenia
\ref{ma2:thm:tangent-descriptions}.
Wektor $(v,w)$ jest normalny wtedy i tylko wtedy, gdy
dla każdego $h$ zachodzi
$\langle v,h\rangle+\langle w,D\varphi(a)h\rangle=0$,
czyli $\langle v+D\varphi(a)^*w,h\rangle=0$.
Podstawiając $h=v+D\varphi(a)^*w$, uzyskujemy
$v=-D\varphi(a)^*w$; odwrotna implikacja wynika
z podstawienia tej równości.
\end{proof}

\section{Immersja a zanurzenie}

Dla otwartego $V\subseteq\R^k$ immersją klasy $C^s$
nazywamy odwzorowanie $\sigma:V\to\R^m$ tej klasy,
którego pochodna jest różnowartościowa w każdym punkcie.
Zanurzenie jest immersją, która jest homeomorfizmem
na swój obraz z topologią podprzestrzeni. Warunek
topologiczny jest istotny także przy różnowartościowej
immersji. Lokalna postać immersji została udowodniona
w twierdzeniu o rzędzie stałym
(\ref{ma2:cor:rank-maximal}).

\begin{stwierdzenie}
\label{ma2:prop:embedding-image}
Jeżeli $V\subseteq\R^k$ jest otwarty, a
$\sigma:V\to\R^m$ jest zanurzeniem klasy $C^s$,
to $S=\sigma(V)$ jest podrozmaitością zanurzoną
wymiaru $k$ tej klasy. Odwrotnie, każda parametryzacja
lokalna podrozmaitości jest zanurzeniem na swój obraz.
\end{stwierdzenie}
\begin{proof}
Ustalmy $a\in V$, $p=\sigma(a)$. Lokalna postać
immersji z wniosku \ref{ma2:cor:rank-maximal} daje
otoczenia $V_0$ punktu $a$ i $U_0$ punktu $p$,
dyfeomorfizmy $\eta,\alpha$ oraz
$\alpha(\sigma(x))=(\eta(x),0)$ dla $x\in V_0$.
Możemy założyć, że $\sigma(V_0)\subseteq U_0$.
Ponieważ $\sigma$ jest homeomorfizmem na $S$,
$\sigma(V_0)$ jest względnie otwarty w $S$;
istnieje zatem otwarty zbiór $O\subseteq\R^m$
taki, że $S\cap O=\sigma(V_0)$. Punkt $p$ należy do $O$.
W otwartym zbiorze $\alpha(U_0\cap O)$ wybieramy
produkt kul $P\times Q$ o środkach w zerze, tak mały,
aby $P\subseteq\eta(V_0)$.
Niech $U=\alpha^{-1}(P\times Q)$.
Każdy punkt $S\cap U$ należy do $\sigma(V_0)$,
więc jego druga współrzędna w mapie $\alpha$ jest zerowa.
Odwrotnie, dla $u\in P$ istnieje $x=\eta^{-1}(u)\in V_0$;
$\alpha(\sigma(x))=(u,0)$, zatem
$\alpha^{-1}(u,0)=\sigma(x)\in S\cap U$.
Ograniczenie $\alpha$ do $U$ jest mapą plasterkową.
Odwrotna część stwierdzenia wynika bezpośrednio
z definicji parametryzacji.
\end{proof}

\begin{przyklad}[Różnowartościowa immersja, która nie jest zanurzeniem]
\label{ma2:ex:injective-nonembedding}
Niech $V=(-\pi/2,\pi)$ i $\tau(t)=(\sin t,\sin 2t)$.
Pochodna $\tau'(t)=(\cos t,2\cos 2t)$ nigdy nie
znika: gdy $\cos t=0$, mamy $\cos 2t=-1$.
Zatem $\tau$ jest immersją.
Odwzorowanie jest różnowartościowe: jeśli $\tau(t)=\tau(s)$,
to $\sin t=\sin s$. Gdy ta wspólna wartość jest
niezerowa, z $\sin 2t=2\sin t\cos t$ otrzymujemy
$\cos t=\cos s$, a zatem $t-s\in2\pi\mathbb Z$.
Długość $V$ jest mniejsza niż $2\pi$, więc $t=s$.
Gdy wspólna wartość jest zerowa, jedyną możliwością
w $V$ jest $t=s=0$.
Jednak dla $t_j=\pi-1/j$ mamy
$\tau(t_j)\to(0,0)=\tau(0)$, podczas gdy $t_j\not\to0$.
Wszystkie $t_j$ dla $j\ge1$ należą do $V$, a
ciąg zbiega w $\R$ do $\pi\ne0$.
Odwrotność na obrazie nie jest ciągła w $(0,0)$,
więc $\tau$ nie jest zanurzeniem.
\end{przyklad}

\begin{przyklad}[Samoprzecięcie immersji]
\label{ma2:ex:immersed-crossing}
Krzywa $\sigma:\R\to\R^2$, $\sigma(t)=(\sin t,\sin 2t)$
jest immersją: $\sigma'(t)=(\cos t,2\cos 2t)$ nie
może być zerem, gdyż $\cos t=0$ pociąga
$\cos 2t=2\cos^2t-1=-1$. Jednak
$\sigma(0)=\sigma(\pi)=(0,0)$.
Prędkości w tych dwóch parametrach są równe
$(1,2)$ i $(-1,2)$; wyznacznik macierzy o tych
kolumnach jest równy $4$, więc są liniowo niezależne.
Gdyby obraz był jednowymiarową podrozmaitością
w otoczeniu zera, obie należałyby do jednowymiarowej
przestrzeni stycznej, co jest niemożliwe.
\end{przyklad}

Dodatkowy warunek właściwości różnowartościowej immersji
(każdy przeciwobraz zbioru zwartego jest zwarty)
zapewnia zanurzenie i domknięty obraz w przypadku
gładkich rozmaitości; dla zwartej dziedziny wystarcza
różnowartościowość immersji. Pełny wynik i dowód znajdują się w
\atlasref{DG-16::prop:proper-immersja}{DG-16: kiedy immersja jest zanurzeniem}.

\section{Przykłady i geometria normalnych}

\begin{przyklad}[Paraboloida]
\label{ma2:ex:submanifold-paraboloid}
Niech $S=\{(x,y,z):z=x^2+y^2\}$.
Jest to wykres gładki, a parametryzacja
$\sigma(x,y)=(x,y,x^2+y^2)$ daje w punkcie
$p=(a,b,a^2+b^2)$, na mocy twierdzenia
\ref{ma2:thm:tangent-descriptions} i wzoru
$D\sigma(a,b)(h,k)=(h,k,2ah+2bk)$, przestrzeń
\[
 T_pS=\{(h,k,2ah+2bk):h,k\in\R\}.
\]
Funkcja $F(x,y,z)=z-x^2-y^2$ ma gradient
$(-2x,-2y,1)$. Jest on niezerowy w każdym punkcie,
a $S=F^{-1}(\{0\})$, więc z twierdzenia
\ref{ma2:thm:tangent-descriptions} otrzymujemy
$N_pS=\operatorname{span}\{(-2a,-2b,1)\}$.
W punkcie $p=0$ przestrzenie są odpowiednio
$\R^2\times\{0\}$ oraz $\{0\}\times\R$.
Płaszczyzna afiniczna styczna w tym punkcie jest
płaszczyzną $z=0$.
\end{przyklad}

\begin{figure}[htbp]
\centering
\begin{tikzpicture}[x={(2cm,-.52cm)},y={(1.2cm,2.1cm)},z={(0cm,.7cm)},>=Stealth,line cap=round,line join=round,font=\small]
\definecolor{manifoldblue}{RGB}{35,91,140}
\definecolor{manifoldred}{RGB}{170,65,55}
\definecolor{manifoldgold}{RGB}{205,145,35}
\filldraw[fill=manifoldgold!15,draw=manifoldgold!80!black,thick] (-1.23,-1.23,0)--(1.23,-1.23,0)--(1.23,1.23,0)--(-1.23,1.23,0)--cycle;
\shade[top color=manifoldblue!12,bottom color=manifoldblue!39]
 plot[domain=-1:1,samples=41,variable=\t] (\t,-1,{\t*\t+1}) --
 plot[domain=-1:1,samples=41,variable=\t] (1,\t,{1+\t*\t}) --
 plot[domain=1:-1,samples=41,variable=\t] (\t,1,{\t*\t+1}) --
 plot[domain=1:-1,samples=41,variable=\t] (-1,\t,{1+\t*\t}) -- cycle;
\draw[manifoldblue!80!black,thick]
 plot[domain=-1:1,samples=41,variable=\t] (\t,-1,{\t*\t+1}) --
 plot[domain=-1:1,samples=41,variable=\t] (1,\t,{1+\t*\t}) --
 plot[domain=1:-1,samples=41,variable=\t] (\t,1,{\t*\t+1}) --
 plot[domain=1:-1,samples=41,variable=\t] (-1,\t,{1+\t*\t}) -- cycle;
\draw[manifoldblue!60,thin] plot[domain=-1:1,samples=41,variable=\t] (\t,0,{\t*\t});
\draw[manifoldblue!60,thin] plot[domain=-1:1,samples=41,variable=\t] (0,\t,{\t*\t});
\draw[->,manifoldred,very thick] (0,0,0)--(0,0,1) node[above left] {$n=(0,0,1)$};
\fill[manifoldred] (0,0,0) circle (2pt) node[below right] {$p$};
\node[manifoldgold!85!black,below] at (1.23,-1.23,0) {$T_0S$};
\node[manifoldblue!80!black,above right] at (1,1,2) {$S$};
\end{tikzpicture}
\caption{Paraboloida, płaszczyzna styczna i kierunek normalny w zerze.}
\label{ma2:fig:submanifold-paraboloid}
\end{figure}

\paragraph{Rysunek interaktywny 3D.}
Rysunek \ref{ma2:fig:submanifold-paraboloid} przedstawia
powierzchnię, jej płaszczyznę styczną i jednostkową
normalną. Wersja Plotly pozwala oglądać tę konstrukcję
z różnych stron. Odcinek normalny jest tylko fragmentem
całej prostej normalnej.
\begin{lstlisting}[language=Python]
import numpy as np
import plotly.graph_objects as go

x = np.linspace(-1, 1, 61)
X, Y = np.meshgrid(x, x)
fig = go.Figure([
    go.Surface(x=X, y=Y, z=X**2+Y**2,
        colorscale=[[0, "#235b8c"], [1, "#d4e5f1"]],
        showscale=False,
        name="S", showlegend=True),
    go.Surface(x=X, y=Y, z=np.zeros_like(X),
        opacity=0.45,
        colorscale=[[0, "#cd9123"], [1, "#f5dfac"]],
        showscale=False, name="T0S", showlegend=True),
    go.Scatter3d(x=[0,0], y=[0,0], z=[0,1],
        mode="lines", name="n", line=dict(width=7,
        color="#aa4137")),
    go.Scatter3d(x=[0], y=[0], z=[0], mode="markers",
        name="p", marker=dict(color="#aa4137", size=5)),
])
fig.update_layout(scene=dict(aspectmode="data",
    xaxis_title="x", yaxis_title="y", zaxis_title="z"))
fig.write_html("am2_21_tangent.html", include_plotlyjs=True)
fig.show()
\end{lstlisting}

\begin{przyklad}[Regularny poziom o kowymiarze dwa]
\label{ma2:ex:submanifold-circle}
Niech $F:\R^3\to\R^2$ będzie dane wzorem
$F(x,y,z)=(x^2+y^2-1,z)$. Na zbiorze $F=0$
mamy $x^2+y^2=1$, więc wektory gradientów
$(2x,2y,0)$ i $(0,0,1)$ są liniowo niezależne.
Pochodna ma rząd $2$, więc stosujemy drugi warunek
twierdzenia \ref{ma2:thm:submanifold-descriptions}
w każdym punkcie zbioru zer. Okrąg jednostkowy
w płaszczyźnie $z=0$ jest podrozmaitością wymiaru $1$.
Dla $p=(a,b,0)$, gdzie $a^2+b^2=1$, mamy
$DF(p)(h,k,l)=(2ah+2bk,l)$, stąd
$\ker DF(p)=\{(h,k,l):2ah+2bk=0,\ l=0\}$.
Z twierdzenia \ref{ma2:thm:tangent-descriptions} dostajemy
\[
 T_pS=\operatorname{span}\{(-b,a,0)\},\qquad
 N_pS=\operatorname{span}\{(a,b,0),(0,0,1)\}.
\]
Pierwszy wektor jest niezerowy i spełnia oba równania,
a dwa wektory normalne są niezależne, więc podane
rozpiętości mają właściwe wymiary.
\end{przyklad}
